%!TEX root = main.tex

\section{Perspectives}
\label{sec:perspectives}

Our experiments show that parameters such as mobility and density strongly impact the performance of broadcast algorithms.
%
Interestingly, these parameters can be monitored at runtime.
%
For instance, nodes can detect the density in their vicinity by sending control messages or by monitoring surrounding traffic.
%
Likewise, mobile devices are often equipped with positioning systems and/or accelerometers to sense the movement of nodes.
%
The next step is using the knowledge of the performance of various protocols in different contexts, to drive an adaptive overlay for disseminating messages.
%
In our vision, a monitoring process provides data to a management layer that decides which approach a node should use to disseminate messages.
%
The main challenge in this case is replacing the algorithm used at runtime without losing any message.
%
Fortunately, most dissemination protocols do not require keeping a complex state on each node.
%
Hence, we think that they are interchangeable while keeping strong guarantees on message delivery.
%
% We may also consider, for instance, an implementation of \emph{ABBA} ``talking'' to another algorithm through a translation layers. %% ER no idea what you mean here.

%\begin{itemize}
%	\item having obstacles as the main source of interference
%	\item the source of a broadcasting session is choosen in a random way among
%	those who are still online. 
%\end{itemize}

%\raziel{which is the impact of having one leader (for example, the one who leads one protest) to broadcast OR using the same node to perform every broadcast session}
